\section{Dynamical localization in the laboratory}
\label{sec:experiments}

Moore, Robinson, Bharucha, Sundaram and Raizen realized the quantum kicked rotor with ultracold sodium
atoms in 1995 \cite{Moore1995}. The atoms move in a standing wave of near-resonant laser light, whose
intensity pattern acts on them as a potential proportional to $\cos(2k_L x)$, with $k_L$ the wave number of
the light. Switching the standing wave on in short pulses with period $T$ turns this into a kicked rotor, with
the angle $\theta = 2k_L x$ and the momentum measured in units of $2\hbar k_L$. The kick strength $\kappa$ plays
the role of our $K$, and the scaled Planck constant $\bar k = 8\omega_r T$, where $\omega_r$ is the atomic
recoil frequency, plays the role of our $\hbar$. Because $\bar k$ depends on the pulse period, the experiment
can tune how quantum the motion is.

In their main run ($\kappa = 11.6$ with a 10\% uncertainty, $\bar k = 2.0$) the momentum spread first grew
diffusively at the classical rate $\kappa^2/2$ and then stopped after a break time of about 8 kicks. From then
on the momentum distribution stayed exponentially localized, checked out to 50 kicks. They quote a
localization length $\xi = \kappa^2/4\bar k^2 \approx 8.3$ in units of $2\hbar k_L$, defined as the $1/e$ point
of the probability distribution, $P \propto \ee^{-|n|/\xi}$. In our convention, with $|\psi_m|^2 \propto
\ee^{-2|m|/\xi}$, this is twice as long, $\kappa^2/2\bar k^2 = D/\hbar^2$, exactly the estimate of
\cref{eq:estimates}.

They also observed quantum resonances, when $\bar k/2$ was a rational multiple of $2\pi$. At odd multiples of
$\pi$ (our antiresonance, $\hbar = 2\pi$) the first two or three kicks showed the predicted sign flip, with the
second kick undoing the first. One difference from our idealized model: the atoms start in a Gaussian wave
packet rather than a single momentum state, so at resonance their distributions stayed narrow instead of
growing quadratically. For a review of chaos experiments with cold atoms see Saif \cite{Saif2005}.

Dynamical localization is an interference effect between Floquet states, so it needs phase coherence.
Anything that randomizes the phase of the wave function between kicks, such as the spontaneous emission of a
photon or noise in the kick strength, would break the interference that stops the growth. In this respect too
it resembles Anderson localization of electrons, which inelastic scattering destroys.
